\section{Results}
\label{sec:results}

We organize the evidence around seven research questions.
\textbf{RQ1:} where is replay memory-compressing? \textbf{RQ2:} does encoder
plasticity improve patient adaptation? \textbf{RQ3:} is the gain caused by
replay, plasticity, or both? \textbf{RQ4:} is the 16\,KiB operating point itself
a limitation? \textbf{RQ5:} does the result survive a class-reweighted
source model? \textbf{RQ6:} does it transfer across databases? \textbf{RQ7:} how
many labels are required?

\paragraph{Which results are pre-specified.} RQ1--RQ3 and the source-domain
change analysis follow a pre-specified plan: the primary metric (per-patient
\macrof{} over present N/S/V classes), the unit of analysis (patient), the
all-DS2 cohort, the seed count, the paired test, the effect size, and the
six-comparison Holm family were all fixed in a written analysis plan before the
final runs, and the six comparisons reported in Table~\ref{tab:pairwise_tests}
are exactly that pre-specified family with no additions. \textbf{One departure
from a strict reading of ``pre-specified'' must be stated:} record 202 was
excluded \emph{after} the subject-identity audit described in
Section~\ref{sec:protocol}, and every affected analysis was rerun. The exclusion
changed the cohort, not the methods --- no metric, test, effect size, or
comparison family was altered in response to seeing results. RQ4--RQ7 are \textbf{secondary and exploratory}: they were specified
after the primary analysis, they use reduced panels or different cohorts, and
their $p$-values are not part of the corrected family. We flag the cohort and
seed count at the head of each secondary subsection because they differ.

\subsection{RQ1 --- The replay-point cliff}

\textbf{Takeaway:} only the post-pool representation is storage-compressive;
every intermediate tap is several times larger than the raw beat, so the
architecture fixes \emph{where} replay can be cheap.

The cliff is a deterministic property of the architecture
(Table~\ref{tab:replay_bytes}, Fig.~\ref{fig:cliff}). A raw beat costs 198\,B,
the convolution/encoder tokens $1{,}056$\,B ($5.3\times$), the FFN hidden tensor
$8{,}448$\,B ($42.7\times$), and only the post-pool$+$RR vector is smaller at
18\,B ($0.09\times$). Consequently, the frozen source model can afford 705
post-pool exemplars but only 62 raw exemplars in the same 16\,KiB
(Table~\ref{tab:arm_definitions}). RQ1 confirms that the cliff constrains
\emph{where} replay is cheap; the remaining questions ask whether cheap replay
is the best use of the budget.

\subsection{RQ2 --- Patient-level adaptation gain}

\textbf{Takeaway:} both encoder-LoRA arms beat the frozen model with a large
effect (significant after Holm correction under the pre-registered primary test;
Table~\ref{tab:pairwise_tests}) and improve 16 of 21 patients, against 12 of 21
for maximum head-only replay, with stronger pooled minority-class performance.

Table~\ref{tab:primary_arms} reports the all-DS2 primary result: per-patient
\macrof{} (seeds median-aggregated within patient, the pre-registered locked
estimator), the number of patients improved versus A0, and pooled S-class F1.
Both encoder-LoRA arms dominate the frozen model and improve the most patients.
A5 reaches the best per-patient score ($0.811\pm0.179$) and A4 nearly matches it
($0.810\pm0.176$) at the same 16/21 improved count; by contrast,
maximum-volume head-only replay (A1) improves only 12/21 patients despite storing
$\approx 11\times$ more exemplars. Fig.~\ref{fig:primary} shows
the per-patient distribution and Fig.~\ref{fig:waterfall} the sorted per-patient
gain; both make clear that the mean is carried by a minority of patients.

\paragraph{Seed stability.} A descriptive point that costs nothing and does not
depend on the inconclusive superiority comparison: the arms differ sharply in how
much their arm-level score depends on the seed-aggregation rule. Comparing
mean-over-seeds and median-over-seeds central values
(\texttt{results/r0\_estimator\_matrix.csv}), maximum head-only replay (A1) moves
$0.786 \to 0.743$, a $0.043$ gap---its arm mean sits well above its median because
a minority of favourable seeds inflate it---while rank-1 LoRA (A4) moves
$0.803 \to 0.777$ ($0.026$) and rank-2 LoRA (A5) is essentially
aggregation-invariant at $0.809$ versus $0.813$ ($0.004$). Maximum head-only
replay is thus the \emph{least seed-stable} arm tested and the proposed encoder
arms the most stable. No test was pre-registered for this contrast, so we report
it descriptively.
Aggregate means hide the fact that adaptation is not uniformly beneficial, so we
report the harm side explicitly, and we report it at two thresholds because the
choice of threshold changes the story materially.

Counting any decline as harm, the encoder arms harm two (A4) and three (A5) of 21
patients and maximum head-only replay harms 8, which makes the encoder arms look
safer. That impression does not survive a meaningful threshold. Defining the harm
rate $R_{\mathrm{harm}} = |\{p : G_p < -\epsilon\}| / P$ with $\epsilon = 0.02$
stated before analysis, the counts become A1 $0/21$
($R_{\mathrm{harm}} = 0.000$), A2 $2/21$, A3 $3/21$, A4 $2/21$, and A5 $2/21$
(Table~\ref{tab:harm}) --- so on this measure \textbf{maximum head-only replay is
the \emph{safest} arm}, harming nobody measurably, while the encoder arms each
harm two patients. \textbf{Most of the apparent harm
advantage of the encoder arms is sub-threshold noise}: at $\epsilon = 0.02$ the
adapted arms harm zero (A1) to three (A3) patients each. The same correction
applies to the improvement counts: the headline 16/21 for A4 and A5 falls to
12/21 when small changes are excluded, against 10/21 for A1. The encoder arms
remain ahead on the improvement count, but by two patients rather than four, and
they are \emph{worse} on the harm count. What separates them is the size of the
tail loss rather than its frequency: A4's worst patient loses $0.029$ where A3's
loses $0.050$ and A5's $0.112$.

We keep the threshold-free counts in Table~\ref{tab:primary_arms} because they
are what the pre-specified plan defined, and report the thresholded counts beside
them so that neither is read alone. Either way the practical conclusion is the
same: personalization helps most patients, measurably harms a few, and a
deployment would need a per-patient acceptance check --- adapt, evaluate on
held-out adaptation-side data, and fall back to the frozen model if the gain is
not positive --- rather than an assumption that adaptation always helps.

\begin{table}[H]
\centering
\small
\caption{Patient-harm profile on the 21-record patient-disjoint cohort, seeds
averaged. $G_p$ is the per-patient change in the primary metric against the
frozen model A0. Improved / unchanged / harmed use a meaningful-change threshold
$\epsilon = 0.02$ stated before analysis, and
$R_{\mathrm{harm}} = |\{p : G_p < -\epsilon\}| / P$. ``Impr.\ ($>0$)'' repeats
the threshold-free count from Table~\ref{tab:primary_arms}.}
\label{tab:harm}
\setlength{\tabcolsep}{4pt}
\begin{tabular}{lrrrrrrrrrr}
\toprule
Arm & Mean $G_p$ & Med.\ $G_p$ & Impr. & Unch. & Harmed & $R_{\mathrm{harm}}$ & Min $G_p$ & 5th pct. & Max $G_p$ & Impr.\ ($>0$) \\
\midrule
A1 & $+0.119$ & $+0.028$ & 11 & 10 & \textbf{0} & \textbf{0.000} & $-0.017$ & $-0.016$ & $+0.793$ & 12 \\
A2 & $+0.124$ & $+0.040$ & 12 & 8  & 1 & 0.048 & $-0.069$ & $-0.020$ & $+0.789$ & 15 \\
A3 & $+0.109$ & $+0.037$ & 11 & 8  & 2 & 0.095 & $-0.052$ & $-0.036$ & $+0.770$ & 13 \\
A4 & $+0.136$ & $+0.035$ & \textbf{12} & 7 & 2 & 0.095 & $\mathbf{-0.034}$ & $-0.027$ & $+0.804$ & \textbf{17} \\
A5 & $+0.143$ & $+0.037$ & \textbf{12} & 7 & 2 & 0.095 & $-0.123$ & $-0.024$ & $+0.804$ & \textbf{17} \\
\bottomrule
\end{tabular}
\end{table}


% GENERATED FILE -- do not edit by hand.
% Regenerate with: python3 scripts/r1_regen_primary_tables.py
\begin{table}[H]
\centering
\small
\caption{Per-patient \macrof{} over the 21-record patient-disjoint cohort, five seeds \textbf{median-aggregated within patient} (the pre-registered locked estimator). Intervals are patient bootstraps ($10{,}000$ resamples).}
\label{tab:primary_arms}
\begin{tabular}{lrrrl}
\toprule
Arm & Mean & SD & Median & 95\% CI \\
\midrule
A0 (frozen) & $0.666$ & $0.245$ & $0.666$ & $[0.563, 0.766]$ \\
A1 (head $+$ max replay) & $0.785$ & $0.177$ & $0.743$ & $[0.712, 0.859]$ \\
A2 (post-pool $r{=}1$) & $0.792$ & $0.179$ & $0.763$ & $[0.717, 0.866]$ \\
A3 (post-pool $r{=}2$) & $0.767$ & $0.180$ & $0.739$ & $[0.693, 0.844]$ \\
A4 (LoRA $r{=}1$) & $0.810$ & $0.176$ & $0.777$ & $[0.737, 0.884]$ \\
A5 (LoRA $r{=}2$) & $0.811$ & $0.179$ & $0.813$ & $[0.737, 0.886]$ \\
\bottomrule
\end{tabular}
\end{table}


\begin{figure}[H]
\centering
\includegraphics[width=0.92\textwidth]{fig_patient_distribution.pdf}
\caption{Per-patient \macrof{} across the 21-record primary cohort (five seeds
averaged within patient, so each point is one patient). Boxes show the median
and interquartile range; grey lines follow a single patient across arms. The
encoder-LoRA arms (A4, A5) shift the distribution upward, but the spread is
wide and heavily patient-dependent: the mean gain is not evenly distributed.}
\label{fig:primary}
\end{figure}

\begin{figure}[H]
\centering
\includegraphics[width=0.95\textwidth]{fig_patient_waterfall.pdf}
\caption{Per-patient gain against the frozen model, sorted, for maximum
head-only replay (A1) and rank-1 encoder LoRA (A4) on the primary cohort. The
grey band is the meaningful-change threshold $|\Delta| \le 0.02$ \macrof{};
bars inside it are indistinguishable from no change, and red bars are harmed
patients. A small number of patients contribute most of the mean gain. At
$\epsilon = 0.02$ the encoder arm (A4) harms two patients while maximum
head-only replay (A1) harms none.}
\label{fig:waterfall}
\end{figure}

Pooled beat-level behavior tells the same story on minority classes
(Table~\ref{tab:pooled}, Fig.~\ref{fig:class}). The frozen model labels the
overwhelming majority of pooled S beats as normal (S F1 $=0.145$); every
adaptation arm recovers S, but the encoder arms are strongest, reaching S F1
$0.832$--$0.835$ and lifting fusion (F) F1 from $0.000$ to $0.578$--$0.604$ where
the post-pool arms stall near $0.39$. Pooled accuracy rises from $0.873$ (A0) to
$0.974$ (A5).

\begin{table}[H]
\centering
\small
\caption{Pooled beat-level performance over the 21-record primary cohort
($32{,}768$\,N / $1{,}385$\,S / $2{,}564$\,V / $291$\,F test beats). Each
record contributes its seed-mean confusion matrix, so the supports are test
beats rather than seed-multiplied decisions. Best in each column in bold.}
\label{tab:pooled}
\begin{tabular}{lrrrr}
\toprule
Arm & Pooled acc. & S F1 & V F1 & F F1 \\
\midrule
A0 & 0.873 & 0.145 & 0.728 & 0.000 \\
A1 & 0.962 & 0.776 & 0.915 & 0.416 \\
A2 & 0.960 & 0.786 & 0.904 & 0.395 \\
A3 & 0.960 & 0.773 & 0.909 & 0.392 \\
A4 & 0.973 & 0.832 & 0.934 & 0.578 \\
A5 & \textbf{0.974} & \textbf{0.835} & \textbf{0.938} & \textbf{0.604} \\
\bottomrule
\end{tabular}
\end{table}

\begin{figure}[H]
\centering
\includegraphics[width=0.90\textwidth]{fig_class_specific.pdf}
\caption{Class-specific pooled performance with the frozen baseline. Adaptation
concentrates its benefit on the minority S class while preserving N and V;
encoder arms give the largest minority-class gains. The fusion (F) class is not
shown here because it is absent from most patients' test segments; pooled F1
for F is reported in Table~\ref{tab:pooled}, where the encoder arms show their
largest relative improvement ($0.000 \to 0.58$--$0.60$ against $\approx 0.39$
for the post-pool arms).}
\label{fig:class}
\end{figure}

\paragraph{Direct paired tests.} Table~\ref{tab:pairwise_tests} and
Fig.~\ref{fig:forest} report the six pre-registered comparisons under the primary
median-seed aggregation, with Holm correction, superiority, and the
pre-registered equivalence test ($\delta = 0.02$ \macrof{}, 90\% BCa interval of
the median paired difference; Section~\ref{sec:protocol}). Both encoder arms
significantly beat the frozen model: A4--A0 median $+0.040$ (95\% BCa CI
$[0.000,0.137]$, Holm $p=0.007$, rank-biserial $0.91$) and A5--A0 median $+0.038$
(CI $[0.000,0.140]$, Holm $p=0.017$, rank-biserial $0.81$). \textbf{These are the
only two comparisons that survive Holm correction.} The \emph{direct} comparisons
against maximum head-only replay are \emph{inconclusive}, not null: A4--A1 has a
median difference of $+0.0007$ and A5--A1 of $+0.0034$ (both Holm $p=0.079$), yet
their 90\% equivalence intervals ($[0.000,0.024]$ and $[0.000,0.029]$) are not
contained in $[-0.02,+0.02]$, so equivalence within the pre-specified margin is
\emph{not} declared either. This 21-patient design detects a median shift no
smaller than its minimum detectable effect ($\approx 0.010$--$0.040$ \macrof{}
across the family; Table~\ref{tab:pairwise_tests}), and the observed A4--A1
median lies below it. The one head-only contrast that \emph{does} clear the
equivalence test is A4 versus the rank-1-equivalent post-pool adapter A2 (90\% CI
$[0.000,0.016] \subset [-0.02,+0.02]$). The honest reading is a bounded
replay--plasticity \emph{allocation} frontier: encoder plasticity reaches
significance against the frozen baseline where post-pool volume does not and
improves more patients (16/21 vs 12/21), while its direct superiority over
maximum replay is neither established nor excluded at this sample size---not
``depth beats volume.'' \textbf{The A4--A1 equivalence verdict is not robust and
turns on a single record.} With all 21 records the $90\%$ BCa upper bound is
$0.0240$, above the $\delta = 0.020$ margin, so equivalence is \emph{not}
declared; dropping record 232 alone lowers that bound to $0.0198$---below the
margin by just $0.0002$---and flips the verdict to equivalent
(\texttt{results/e8\_sensitivity\_record232.csv}). Record 232 is the
sixth-largest A4$-$A1 gain ($+0.0428$), so the equivalence conclusion rests on one
patient and a $0.0002$-\macrof{} margin and is not robust in either direction. We
report this rather than pick whichever cohort gives the tidier verdict.

\paragraph{Why the paired statistic looks smaller than the arm gap.} A reader
comparing the arm summaries (A4 at $0.810$, A1 at $0.785$) sees a gap of about
$0.025$, whereas the pre-registered test statistic---the \emph{median of the
per-patient paired differences}---is $+0.0007$, roughly $36\times$ smaller. This
is not an arithmetic error: the median of differences equals the difference of
medians only when the difference distribution is symmetric, and here it is
strongly right-skewed. The 21 per-patient A4$-$A1 differences run from $-0.033$ to
$+0.167$ with 14 positive, 3 exactly zero, and 4 negative; three patients alone
hold $61\%$ of the total positive mass. The mean paired difference ($+0.025$) and
the difference of arm central values ($+0.025$ for the locked estimator, $+0.034$
for a median-of-medians summary) are therefore an order of magnitude above the
median paired difference. The pre-registered statistic is the latter, which is why
the direct A4$-$A1 comparison is inconclusive despite the visibly larger arm gap.
The bootstrap interval reflects the same structure: the $95\%$ and $90\%$ BCa
lower bounds are \emph{exactly} $0.000$ because 3 patients have exactly zero
difference and only 4 are negative, so a resampled median can never fall below
zero. This is a property of a skewed, zero-heavy difference distribution, not a
clipping artifact in the estimator.

% Auto-generated by scripts/run_equivalence_mde.py (T2). Do not hand-edit.
\begin{tabular}{lrrrrrc}
\toprule
Comparison & Median $\Delta$ & 95\% BCa CI & Holm $p$ & $r_{rb}$ & 90\% CI (TOST) & Equiv.\ ($\delta{=}0.02$) \\
\midrule
A4--A1 & +0.0007 & [+0.000, +0.024] & 0.079 & +0.72 & [+0.000, +0.024] & -- \\
A5--A1 & +0.0034 & [+0.000, +0.029] & 0.079 & +0.74 & [+0.000, +0.029] & -- \\
A4--A2 & +0.0044 & [+0.000, +0.016] & 0.079 & +0.63 & [+0.000, +0.016] & \checkmark \\
A4--A3 & +0.0043 & [+0.000, +0.037] & 0.079 & +0.67 & [+0.000, +0.037] & -- \\
A4--A0 & +0.0401 & [+0.000, +0.137] & 0.007 & +0.91 & [+0.001, +0.137] & -- \\
A5--A0 & +0.0377 & [+0.000, +0.140] & 0.017 & +0.81 & [+0.001, +0.140] & -- \\
\bottomrule
\end{tabular}


\begin{figure}[H]
\centering
\includegraphics[width=0.82\textwidth]{fig05_paired_differences.pdf}
\caption{Planned paired comparisons: median paired \macrof{} difference with the
90\% BCa interval, against the pre-registered $\pm\delta$ equivalence band
($\delta = 0.02$). Encoder arms clear zero against the frozen model (A4--A0,
A5--A0); against maximum head-only replay (A4--A1, A5--A1) the intervals are
neither below the significance threshold after correction nor inside the
equivalence band---bounded but inconclusive.}
\label{fig:forest}
\end{figure}

\subsection{RQ3 --- Controlled replay-versus-plasticity ablations}

\textbf{Takeaway:} the encoder-plasticity contrast is statistically supported
(B4$-$B2, $p=0.005$, 16 of 21 patients improved), while replay \emph{volume} is
not; replay's clearest measured role is source retention, not target gain.

To separate replay location, replay volume, and trainable depth, we run the
B0--B7 controlled ablation matrix (Table~\ref{tab:causal_ablation},
Fig.~\ref{fig:causal}).
Arm means suggest a tidy story: adding raw replay to encoder LoRA lifts mean
\macrof{} from $0.755$ (B3, LoRA, no replay) to $0.803$ (B4, LoRA $+$ 62 raw), a
$+0.048$ apparent contribution of \emph{replay beyond plasticity}, while moving
from a frozen-encoder head at the same raw-replay budget (B2) to encoder LoRA
(B4) adds $+0.030$ from \emph{plasticity beyond replay}.

\paragraph{Only one of those two contributions survives paired testing.}
Differences of arm means are not evidence of a per-patient effect, so we test
each contrast patient-wise (Table~\ref{tab:causal_contrasts}). The plasticity
contrast is solid: B4$-$B2 gives $+0.030$ with a 95\% interval of
$[+0.007,+0.063]$ excluding zero, Wilcoxon $p = 0.005$, rank-biserial $0.74$, and
16 of 21 patients improved against 4 worsened. \textbf{The replay contrast is
not.} B4$-$B3 has a mean of $+0.048$ but a \emph{median of exactly zero}, an
interval $[-0.006,+0.112]$ that includes zero, $p = 0.473$, and a 10/11 split
between patients helped and harmed --- marginally more patients are \emph{worse}
with replay added than better. Its mean is carried by a handful of patients with
large gains, not by a consistent effect.

The two remaining contrasts sharpen the paper's central point. Replay
\emph{volume} does almost nothing: B1$-$B6 compares 705 post-pool exemplars
against 62 and yields $+0.008$ ($p = 0.232$) with only 5 patients improved and
\textbf{15 worsened} --- spending an extra 643 exemplars makes three times as
many patients worse as better. Encoder depth at matched post-pool volume
(B4$-$B7) is $+0.013$ and also not significant ($p = 0.179$).

We therefore state the controlled comparison result as: \emph{controlled ablations show a
statistically supported contribution from encoder plasticity, and a positive but
statistically unsupported mean contribution from replay, with no measurable
benefit from replay volume.} The earlier framing that replay and plasticity
``contribute independently'' overstated what the ablations establish, and we do
not use it. Replay's clearest measured role in this paper is not target gain at
all --- it is source retention, where its effect is large and consistent
(Section~\ref{sec:reg} and Section~\ref{sec:srcchange}).

\begin{table}[H]
\centering
\small
\caption{Paired uncertainty for the four controlled contrasts on the 21-record
patient-disjoint cohort, seeds averaged within patient. Arm means alone are
misleading here: only the \emph{plasticity} contrast is statistically supported.
``Mean''/``Med.'' are the paired difference, ``95\% CI'' a patient bootstrap
interval over the 21 patients ($10{,}000$ resamples), $p$ the Wilcoxon signed-rank value,
$r_{\mathrm{rb}}$ the rank-biserial effect size, ``I/W'' the threshold-free
improved/worsened split, and ``@$\epsilon$'' the same split at the $0.02$
threshold. B1, B4, and B7 are configuration-identical to A1, A4, and A2.}
\label{tab:causal_contrasts}
\setlength{\tabcolsep}{3pt}
\begin{tabular}{llrrcrrrr}
\toprule
Contrast & Interpretation & Mean & Med. & 95\% CI & $p$ & $r_{\mathrm{rb}}$ & I/W & @$\epsilon$ \\
\midrule
B4$-$B3 & replay beyond plast.      & $+0.048$ & $-0.000$ & $[-0.006,+0.112]$ & $0.473$ & $0.15$ & 10/11 & 6/4 \\
\textbf{B4$-$B2} & \textbf{plast.\ beyond replay} & $+0.030$ & $+0.006$ & $\mathbf{[+0.007,+0.063]}$ & $\mathbf{0.005}$ & $\mathbf{0.74}$ & \textbf{16/4} & 6/0 \\
B1$-$B6 & replay volume             & $+0.008$ & $-0.002$ & $[-0.027,+0.056]$ & $0.232$ & $0.37$ & 5/15 & 4/6 \\
B4$-$B7 & encoder depth             & $+0.013$ & $+0.003$ & $[-0.019,+0.043]$ & $0.179$ & $0.40$ & 14/6 & 8/4 \\
\bottomrule
\end{tabular}
\end{table}


\begin{table}[t]
\centering
\small
\caption{E9 controlled replay-versus-plasticity ablation matrix, recomputed on the
21-record patient-disjoint cohort (five seeds). \textbf{Arm means only} --- the
paired uncertainty that decides which contrasts are real is in
Table~\ref{tab:causal_contrasts}, and the mean differences here must not be read
as effects.}
\label{tab:causal_ablation}
\begin{tabular}{lllllll}
\toprule
ID & Status & Scope & Replay & Items & Mean & vs B1 \\
\midrule
B0 & complete & Frozen & none & 0 & 0.666 & $-0.119$ \\
B1 & complete & Head & postpool & 705 & 0.786 & $0.000$ \\
B2 & complete & Head & raw & 72 & 0.773 & $-0.013$ \\
B3 & complete & Encoder LoRA rank 1 & none & 0 & 0.755 & $-0.031$ \\
B4 & complete & Encoder LoRA rank 1 & raw & 62 & 0.803 & $+0.017$ \\
B5 & infeasible & Encoder LoRA rank 1 & raw & 72 &  &  \\
B6 & complete & Head & postpool & 62 & 0.778 & $-0.008$ \\
B7 & complete & Post-pool adapter rank 1 & postpool & 678 & 0.790 & $+0.004$ \\
\bottomrule
\end{tabular}
\end{table}


\begin{figure}[H]
\centering
\includegraphics[width=0.86\textwidth]{fig_causal_ablation.pdf}
\caption{Replay-versus-plasticity ablation, arm means only. The apparent
ordering B4$>$B3 (replay added) and B4$>$B2 (plasticity added) is what mean
differences suggest, but \textbf{only the plasticity contrast survives paired
testing}: B4$-$B2 has $p = 0.005$ while B4$-$B3 has $p = 0.473$ with a median of
essentially zero and a 10/11 patient split. Read this figure together with
Table~\ref{tab:causal_contrasts}, which carries the uncertainty; the bars alone
overstate the replay effect.}
\label{fig:causal}
\end{figure}

\subsection{Can regularization substitute for replay?}
\label{sec:reg}

\textbf{Takeaway:} replay provides substantially better source retention than the
tested zero-byte $B$-factor anchor under the evaluated configuration --- at
matched bytes the anchor loses about five times as much source performance as
replay while adapting less. One regularizer was tested, so this is not a
statement about regularization in general.

\paragraph{What this does and does not say about EWC-LoRA.} The closest
published method is EWC-LoRA \cite{zheng2026ewclora}, which applies a
Fisher-based penalty to a shared low-rank update for task-sequential continual
learning. \textbf{We do not implement it}, so nothing here should be read as a
comparison against it. Two differences make a direct transfer non-obvious rather
than merely untested: their setting accumulates \emph{tasks} while ours is
constrained by \emph{bytes}, and a Fisher importance estimate is persistent
state that would itself have to be charged against the 16\,KiB ceiling --- a
diagonal Fisher over even the 223 trainable parameters of A4 is small, but over
the full 6{,}643-parameter model it is not (Section~\ref{sec:relatedwork}).
Evaluating EWC-LoRA under a byte budget is the single most relevant extension of
the regularization arm of this study, and we flag it as such rather than
claiming the question is closed.

The arms compared so far are all replay-based, which leaves open whether the
budget would be better spent on a regularizer than on stored exemplars. The
classical answer, elastic weight consolidation, is infeasible here: a full-model
FP32 parameter anchor plus a Fisher value costs $6643 \times 8 = 53{,}144$\,B,
more than three times the entire budget before a single exemplar is stored.

A low-rank adapter admits a much cheaper anchor. The attention delta is $(xA)B$
with $B$ zero-initialised, so shrinking $B$ pulls the adapted model back toward
the frozen source function, and it costs \textbf{zero persistent bytes} because
the anchor \emph{is} the frozen base weight the device already stores for
inference. The implemented penalty is
\begin{equation}
\mathcal{L}_{\mathrm{reg}} = \lambda \sum_{m \in \{q,k,v,o\}} \|B_m\|_F^2,
\qquad \lambda = 10^{-2},
\label{eq:anchor}
\end{equation}
applied to the $B$ factor of each of the four attention projections. We state one
caveat plainly, because it bounds the interpretation: the quantity that actually
determines the function is the product $A_mB_m$, and since
$A_mB_m = (cA_m)(B_m/c)$, penalising $B_m$ alone constrains the
\emph{parametrisation} rather than the function. $\|A_mB_m\|_F^2$ is the
better-posed penalty; Eq.~\eqref{eq:anchor} is what was run, and we report it as
such rather than describing it as a direct penalty on the source-function
deviation. This is therefore \textbf{a low-cost memory-matched regularization
baseline}, not the strongest regularizer the budget admits --- restricted EWC on
the adapter parameters, source-logit distillation, and output-consistency
regularization all remain untested. We add two arms: R1 (rank-1 LoRA, anchor, no
replay), byte-comparable with B3, and R2 (rank-1 LoRA, anchor, 62 raw exemplars),
byte-identical to A4.

The answer is clear in both directions (Table~\ref{tab:regularization}).
\emph{As a substitute for replay, the anchor is weak.} Without replay it trades a
little target gain for a little stability relative to unregularized LoRA
($0.755 \to 0.744$ target, $\Delta^{\mathrm{src}}$ $-0.159 \to -0.136$). Replay achieves far
more on both axes at once: A4 reaches $+0.136$ target gain against R1's $+0.078$,
and holds the source-domain change to $-0.027$ against R1's $-0.136$, so the anchor loses
about five times as much source performance as replay does, while also adapting
less. \emph{As a
complement to replay, the anchor is redundant.} R2 is statistically
indistinguishable from A4 on target performance ($0.803$ versus $0.803$) and
improves $\Delta^{\mathrm{src}}$ only from $-0.027$ to $-0.023$.

We take two things from this, both hedged by the caveat above. First, the paper's
replay-centric framing is not an untested assumption: at this scale, stored source
exemplars regulate stability substantially better than this parameter-space
penalty does, which is what one would expect when the shift is a change of patient
rather than a change of task. Second, the anchor is nearly free, so R2 remains a
reasonable default for a deployment wanting the small extra retention margin at no
byte cost. Because only one regularizer was tested, and in its scale-ambiguous
form, we claim only that \emph{this} anchor does not substitute for replay --- not
that regularization in general cannot.

\begin{table}[H]
\centering
\small
\caption{Memory-matched regularization baseline on the 21-record patient-disjoint
cohort (5 seeds). The source anchor is $\lambda\sum_m\|B_m\|_F^2$ on the $B$
factor of each attention projection (Eq.~\ref{eq:anchor}); because $B$ is
zero-initialised and the frozen base weights \emph{are} the anchor, it costs
\textbf{zero persistent bytes}, so R1 is byte-comparable with B3 and R2 with A4.
Penalising $B$ rather than the product $AB$ is scale-ambiguous and constrains the
parametrisation rather than the function; this is a low-cost baseline, not the
strongest regularizer available. $\Delta^{\mathrm{src}}$ is the signed
source-domain \macrof{} change.}
\label{tab:regularization}
\setlength{\tabcolsep}{4pt}
\begin{tabular}{llccrrrr}
\toprule
Arm & Mechanism & Replay & Anchor & Bytes & \macrof{} & $\Delta$ vs.\ A0 & $\Delta^{\mathrm{src}}$ \\
\midrule
A0 & frozen                    & ---  & ---  & 0        & 0.666 & ---                        & $0.000$ \\
B3 & LoRA $r{=}1$              & ---  & ---  & 3{,}568  & 0.755 & $+0.089$ ($p=0.044$)       & $-0.159$ \\
R1 & LoRA $r{=}1$ $+$ anchor   & ---  & yes  & 3{,}568  & 0.744 & $+0.078$ ($p=0.052$)       & $-0.136$ \\
A4 & LoRA $r{=}1$ $+$ replay   & 62   & ---  & 16{,}208 & 0.803 & $+0.136$ ($p=0.0012$)      & $\mathbf{-0.027}$ \\
R2 & LoRA $r{=}1$ $+$ both     & 62   & yes  & 16{,}208 & \textbf{0.803} & $\mathbf{+0.137}$ ($p=0.0011$) & $-0.023$ \\
\bottomrule
\end{tabular}
\end{table}


\subsection{RQ4 --- Adam budget sweep for rank-1 and rank-2 adaptation}
\label{sec:rq4}

\textbf{Takeaway:} within the measured 12-record Adam rank-1/rank-2 panel,
raising the adaptation-state budget from 8 to 64\,KiB produced no statistically
detectable benefit (secondary evidence, reduced panel). We do not generalize
beyond the panel that was measured.

The primary study fixes a single 16\,KiB budget, which raises the obvious
question of whether the conclusions are an artifact of that particular point. The
experiment that speaks to it, E10, was designed as a full budget $\times$
optimizer $\times$ rank grid but was executed only in part, so we name it for
what was actually measured: an \textbf{Adam budget sweep for rank-1 and rank-2
adaptation} at 8, 16, 32, and 64\,KiB. The unmeasured dimensions are not
discussed as results anywhere in this paper.

\paragraph{Scope, stated up front.} This sweep is \textbf{secondary evidence and
is not comparable to the primary cohort}. Two limitations must be read before the
numbers. First, to keep the grid tractable it was executed on a reduced panel of
\textbf{12 DS2 records and 3 seeds} (864 adaptation cells), not the 21-record,
5-seed primary cohort; absolute values are therefore not interchangeable with
Table~\ref{tab:primary_arms}, and no cell of this sweep may be ranked against a
primary-cohort arm. Second, the executed grid is a \emph{subset} of the planned
design: 49 of 140 planned configurations are measured, 5 are analytically
infeasible at their budget, and 86 remain unrun. The budget $\times$ scope
$\times$ rank surface is complete \emph{only for Adam at ranks 1 and 2}; SGD is
measured at 8\,KiB alone, SGD-with-momentum not at all, and ranks 4 and 8 are
almost entirely unmeasured (rank 8 is infeasible at 8 and 16\,KiB by the memory
model).

The single claim this experiment supports is therefore narrow and we state it
once, here: \emph{within the measured Adam rank-1/rank-2 panel, no statistically
resolvable improvement was observed when increasing the persistent-state budget
from 8 to 64\,KiB.} We make \textbf{no optimizer-ranking claim, no rank-4 or
rank-8 claim, and no claim that the 140-cell design is complete.}

\paragraph{Result: no detectable benefit above 8\,KiB.} Within the Adam panel
(Table~\ref{tab:e10_adam}, Fig.~\ref{fig:saturation}), raising the budget from 8
to 64\,KiB does not improve per-patient \macrof{} for any family. The encoder-LoRA
family is flat-to-declining ($0.805 \to 0.790 \to 0.776 \to 0.768$), the
post-pool-adapter family varies within $0.762$--$0.780$, and the head-only family
within $0.762$--$0.774$.

Overlapping confidence intervals are not evidence of equivalence, so we test the
budget effect directly. Holding optimizer, trainable scope, rank, and replay
location fixed and varying \emph{only} the budget gives a paired contrast on each
record; we test each with a Wilcoxon signed-rank test and, for equivalence, with
a two-one-sided-tests (TOST) procedure against a margin of $\delta = 0.05$
\macrof{} stated before analysis (Table~\ref{tab:e10_paired}). Of the 18
budget contrasts, \textbf{14 met the pre-specified $\pm 0.05$ equivalence
criterion} against their 8\,KiB reference, and none of the raises is
significantly positive.

\paragraph{Why $\delta = 0.05$, and what the 18 decisions are not.} The margin
was fixed before analysis at $0.05$ \macrof{} because it is roughly a third of
the primary adaptation effect this paper is built on (A4$-$A0 $= +0.136$): a
budget change that moves \macrof{} by less than $\delta$ cannot account for the
adaptation gain, which is the question the sweep exists to answer. It is a
within-panel yardstick, not a claim about clinical meaningfulness, and it is
coarser than the $\epsilon = 0.02$ per-patient harm threshold used in RQ2
because the sweep panel has 12 records rather than 21 and correspondingly wider
intervals. \textbf{The 18 TOST decisions are not multiplicity-adjusted.} They
sit outside the six-comparison Holm family, which was reserved for the primary
analysis, so each equivalence decision is an uncorrected per-contrast statement
and the count ``14 of 18'' should be read descriptively. Adjusting would make
equivalence harder to declare, not easier, so the direction of the resulting
bias is toward over-claiming equivalence --- another reason we keep the claim
inside the measured panel.

Two contrasts deserve explicit mention because they cut against a naive
``more memory is at least harmless'' reading. For rank-2 encoder LoRA, moving
from 8 to 64\,KiB \emph{significantly reduces} \macrof{} by $0.039$
(Wilcoxon $p = 0.003$), and the 32\,KiB contrast is also negative
($-0.027$, $p = 0.027$); neither is equivalent within the margin. The likely
mechanism is that a raw-replay buffer of several hundred source exemplars starts
to dominate a 500-beat target set, pulling the adapted model back toward the
source. We report this as a measured effect on the secondary panel rather than a
general law.

The defensible reading is a null result of a useful kind, stated at the scope it
was measured: \emph{within the measured 12-record Adam rank-1/rank-2 panel,
increasing the adaptation-state budget produced no statistically detectable
benefit, and for rank-2 encoder LoRA it produced a detectable cost}. On that
panel, extra persistent memory did not buy accuracy, which is consistent with
the allocation question studied in RQ2 and RQ3 --- how to spend a fixed budget
--- being the one that matters. We do not extend this to unmeasured optimizers,
ranks, or cohorts, and we do not claim that 16\,KiB is provably sufficient in
general.

\paragraph{A separate low-byte measurement.}\label{sec:lowbyte} Independently of this sweep, a
fixed-scope run on the \emph{full} 22-record, 5-seed cohort measured a rank-1
encoder-LoRA configuration holding only \textbf{8 raw replay samples} in
$5{,}246$\,B --- under a third of the budget --- reaching $0.801$ per-patient
\macrof{} on the contaminated 22-record cohort, against $0.802$ for the 16\,KiB
rank-2 arm (A5) on the same cohort. We quote this pair on the 22-record basis
because that fixed-scope run was not repeated after the cohort correction.
\textbf{It is therefore a historical sensitivity result and is \emph{not}
directly comparable with Table~\ref{tab:primary_arms}}, which is computed on the
corrected 21-record cohort; the two numbers quoted here are internally
comparable with each other and with nothing else in this paper. Read only within
itself, the pair is suggestive of the same saturation direction --- a
third-of-budget configuration landing beside a full-budget one --- but it carries
no weight against any primary-cohort figure.

\begin{table}[H]
\centering
\small
\caption{Adam budget sweep, rank-1 and rank-2 (E10): best configuration per family at each budget
(per-patient \macrof{}, N/S/V present classes). \textbf{Reduced secondary panel
of 12 DS2 records and 3 seeds, Adam only, ranks 1 and 2 --- not comparable to the
primary cohort of Table~\ref{tab:primary_arms}.} SGD and SGD-with-momentum cells
are excluded as unmeasured (except 8\,KiB SGD). Formal budget contrasts and
equivalence decisions are in Table~\ref{tab:e10_paired}; the means here are
descriptive and should not be read as a ranking.}
\label{tab:e10_adam}
\begin{tabular}{lrrrr}
\toprule
Family & 8\,KiB & 16\,KiB & 32\,KiB & 64\,KiB \\
\midrule
Encoder LoRA $+$ raw replay        & \textbf{0.805} & \textbf{0.790} & 0.776 & \textbf{0.768} \\
Post-pool adapter $+$ post-pool    & 0.762 & 0.765 & \textbf{0.780} & 0.767 \\
Head only                          & 0.762 & 0.774 & 0.765 & 0.763 \\
\bottomrule
\end{tabular}
\end{table}

% GENERATED FILE -- do not edit by hand.
% Regenerate with: make tables  (scripts/make_tables.py)
\begin{table}[H]
\centering
\small
\caption{Paired budget contrasts within the Adam rank-1/rank-2 sweep (E10). Optimizer, trainable scope, rank, and replay location are held fixed and only the budget varies, so each row is a paired comparison over the same 12 records of the secondary panel (3 seeds, Adam only). $\Delta$ is the mean paired change in per-patient \macrof{} relative to the same configuration at 8\,KiB. ``Equiv.'' is a TOST equivalence decision at the pre-stated margin $\delta = 0.05$. 14 of 18 contrasts are equivalent. Raising the budget is never significantly beneficial; for rank-2 encoder LoRA at 64\,KiB it is significantly harmful.}
\label{tab:e10_paired}
\setlength{\tabcolsep}{5pt}
\begin{tabular}{llcrrcc}
\toprule
Scope & Replay & Rank & Budget & $\Delta$ vs.\ 8\,KiB & Wilcoxon $p$ & Equiv. \\
\midrule
\multirow{3}{*}{Encoder LoRA} & \multirow{3}{*}{raw} & \multirow{3}{*}{1} & 16\,KiB & $+0.015$ & $0.204$ & yes \\
 & & & 32\,KiB & $+0.007$ & $0.424$ & yes \\
 & & & 64\,KiB & $+0.002$ & $0.380$ & yes \\
\midrule
\multirow{3}{*}{Encoder LoRA} & \multirow{3}{*}{raw} & \multirow{3}{*}{2} & 16\,KiB & $-0.013$ & $0.164$ & yes \\
 & & & 32\,KiB & $-0.027$ & $\mathbf{0.027}$ & yes \\
 & & & 64\,KiB & $-0.039$ & $\mathbf{0.003}$ & no \\
\midrule
\multirow{3}{*}{Head only} & \multirow{3}{*}{post-pool} & \multirow{3}{*}{0} & 16\,KiB & $+0.006$ & $0.577$ & yes \\
 & & & 32\,KiB & $+0.004$ & $0.898$ & yes \\
 & & & 64\,KiB & $+0.003$ & $0.831$ & yes \\
\midrule
\multirow{3}{*}{Head only} & \multirow{3}{*}{raw} & \multirow{3}{*}{0} & 16\,KiB & $+0.000$ & $0.519$ & yes \\
 & & & 32\,KiB & $-0.004$ & $0.677$ & yes \\
 & & & 64\,KiB & $+0.010$ & $0.776$ & no \\
\midrule
\multirow{3}{*}{Post-pool adapter} & \multirow{3}{*}{post-pool} & \multirow{3}{*}{1} & 16\,KiB & $+0.016$ & $0.240$ & yes \\
 & & & 32\,KiB & $+0.030$ & $0.519$ & no \\
 & & & 64\,KiB & $+0.017$ & $0.831$ & no \\
\midrule
\multirow{3}{*}{Post-pool adapter} & \multirow{3}{*}{post-pool} & \multirow{3}{*}{2} & 16\,KiB & $-0.001$ & $1.000$ & yes \\
 & & & 32\,KiB & $+0.007$ & $0.519$ & yes \\
 & & & 64\,KiB & $+0.006$ & $1.000$ & yes \\
\bottomrule
\end{tabular}

\vspace{0.4em}
{\small A contrast can be \emph{both} significantly different from zero and statistically equivalent within $\pm 0.05$: rank-2 encoder LoRA at 32\,KiB is one (Wilcoxon $p = 0.027$, TOST $p = 0.027$). Significance asks whether the difference is distinguishable from zero; equivalence asks whether it is small. The two columns answer different questions and must not be read as each other's complement. The 18 decisions are \textbf{not} multiplicity-adjusted.}
\end{table}


\begin{figure}[H]
\centering
\includegraphics[width=0.80\textwidth]{fig_e10_budget_saturation.pdf}
\caption{Per-patient \macrof{} versus persistent-state budget for the Adam
rank-1/rank-2 sweep (E10). \textbf{Secondary cohort: 12 DS2 records, 3 seeds, Adam only, ranks 1 and
2; not directly comparable to Table~\ref{tab:primary_arms}.} Points are the best
configuration per family at each budget and shaded bands are 95\% bootstrap
intervals over records. The bands overlap, but overlap alone is not evidence of
equality --- the supporting claim is the paired TOST analysis of
Table~\ref{tab:e10_paired}, in which 14 of 18 budget contrasts are equivalent
within $\pm 0.05$ \macrof{} and none is significantly positive.}
\label{fig:saturation}
\end{figure}

\subsection{Does the result depend on spending the last kilobyte?}
\label{sec:reserve}

\textbf{Takeaway:} the result does not depend on saturating the budget ---
reserving 1\,KiB and regenerating replay counts leaves every adapted arm
significantly above frozen ($p\le0.010$) with encoder adaptation still among the
top allocations.

Every configuration reported so far fills the 16\,KiB budget almost exactly ---
A2 reaches $16{,}383$ of $16{,}384$ bytes. That is not a deployable margin: real
firmware needs room for a version field, a checksum, and future additions, and a
configuration generated against a ceiling it exactly saturates would break the
moment anything is added. We therefore re-ran the entire A0--A5 grid, same 21
records and five seeds, against an effective ceiling of
$B_{\mathrm{budget}} - B_{\mathrm{reserve}} = 16{,}384 - 1{,}024 = 15{,}360$\,B,
regenerating replay counts rather than reporting old results against new byte
figures. The reserve costs A1 49 exemplars, A4 five, and A5 five
(Table~\ref{tab:arm_definitions}).

The qualitative conclusion is preserved (Table~\ref{tab:reserve}): every adapted
arm remains significantly above the frozen model ($p \le 0.010$), encoder
adaptation remains among the top-performing allocations, and A4 is unchanged at
$0.802$. The \emph{exact} arm ordering is not preserved, and we say so rather
than claiming otherwise: without the reserve the order is
$\mathrm{A5} > \mathrm{A4} > \mathrm{A2} > \mathrm{A1} > \mathrm{A3}$, and with it
$\mathrm{A4} > \mathrm{A1} > \mathrm{A5} > \mathrm{A2} > \mathrm{A3}$. A1 gains
slightly ($0.786 \to 0.793$) while A5 loses more ($0.809 \to 0.792$). A5 is thus
more sensitive than A4 to losing the last few replay samples, which supports
\textbf{rank-1 encoder LoRA as the more robust engineering operating point} --- a
practical conclusion the unreserved analysis alone would have missed. We keep the unreserved configuration as the primary analysis because it
is what the pre-specified plan defined, and report the reserve replication as the
deployable-margin check. \textbf{The primary statistical conclusion is
unchanged, although the exact ordering and the preferred engineering operating
point change under the reserve.}

\begin{table}[H]
\centering
\small
\caption{Implementation-reserve replication on the 21-record patient-disjoint
cohort. The full A0--A5 grid re-run against an effective ceiling of
$15{,}360$\,B, i.e.\ holding back a 1\,KiB firmware reserve, same records and
five seeds. Replay counts fall by roughly $7$--$8\%$. Every arm still
significantly beats the frozen model and encoder adaptation stays among the top
allocations, though the exact ordering changes: A5 loses more than A4, making
rank-1 the more robust operating point. \emph{Values are seeds mean-averaged
within patient; median unification of this secondary replication is deferred
(see \texttt{BLOCKERS.md}, ESC-4).}}
\label{tab:reserve}
\setlength{\tabcolsep}{5pt}
\begin{tabular}{lrrrrrr}
\toprule
& \multicolumn{3}{c}{16{,}384\,B ceiling} & \multicolumn{3}{c}{$15{,}360$\,B ceiling ($+1$\,KiB reserve)} \\
\cmidrule(lr){2-4}\cmidrule(lr){5-7}
Arm & Items & \macrof{} & $\Delta$ vs.\ A0 & Items & \macrof{} & $\Delta$ vs.\ A0 \\
\midrule
A0 & 0   & 0.666 & ---      & 0   & 0.666 & --- \\
A1 & 705 & 0.786 & $+0.119$ & 656 & 0.793 & $+0.127$ \\
A2 & 678 & 0.790 & $+0.124$ & 629 & 0.779 & $+0.113$ \\
A3 & 650 & 0.775 & $+0.109$ & 601 & 0.771 & $+0.105$ \\
A4 & 62  & 0.803 & $+0.136$ & 57  & \textbf{0.802} & $\mathbf{+0.136}$ \\
A5 & 52  & \textbf{0.809} & $\mathbf{+0.143}$ & 47 & 0.792 & $+0.125$ \\
\bottomrule
\end{tabular}

\vspace{0.4em}
{\small All reserve-arm gains over A0 remain significant by paired Wilcoxon
($p = 0.006$, $0.003$, $0.010$, $0.004$, $0.005$ for A1--A5). Unreserved the
order is $\mathrm{A5}>\mathrm{A4}>\mathrm{A2}>\mathrm{A1}>\mathrm{A3}$; reserved
it is $\mathrm{A4}>\mathrm{A1}>\mathrm{A5}>\mathrm{A2}>\mathrm{A3}$.}
\end{table}


\subsection{Does the boundary filter dependency matter?}
\label{sec:splitfirst}

\textbf{Takeaway:} no. Removing the dependency entirely moves every arm mean by
at most $0.007$ \macrof{}, leaves both encoder-versus-frozen contrasts
significant at $p = 0.0012$, and leaves both encoder-versus-A1 contrasts
non-significant. The conclusions are not artifacts of non-causal preprocessing.

The released preprocessing filters each record as a whole before the
chronological split (Section~\ref{sec:protocol}), so adaptation-side samples
within the filters' backward reach of the boundary can depend on test-segment raw
samples. To remove that dependency rather than argue about its size, we rebuilt
the beat store \emph{split-first}: the raw signal is cut at the boundary, a
$720$-sample ($2$\,s) guard is discarded from each side --- roughly a sixfold
margin on the measured $120$-sample influence reach --- and the two pieces are
filtered independently, so no adaptation-side sample can see a test-side raw
sample. Beats whose window would straddle the cut are dropped, which costs one to
three adaptation beats per record; we then truncate every record to a uniform
$497$-beat adaptation block so the protocol stays balanced.

We reran A0, A1, A4, and A5 on this store over the same 21 patient-disjoint
records and the same five seeds. Table~\ref{tab:splitfirst} reports the result.

Three things are worth separating. First, the \textbf{frozen arm is unmoved}
($-0.0000$, largest per-record shift $0.0004$). A0 performs no adaptation, so it
isolates the effect of refiltering the test segment alone, and that effect is
numerically negligible --- which is the expected result and a useful check that
the rebuild did not perturb the evaluation data.

Second, the \textbf{adapted arms move very little, and not in the direction
leakage would predict}. A1 rises by $0.007$, A4 falls by $0.006$, and A5 is
unchanged to three decimals. If whole-record filtering had been inflating the
results, removing it should have lowered them systematically; instead the shifts
have mixed sign. \textbf{The seeds are identical across the two runs, so these
are not run-to-run fluctuations.} They reflect the combined effects of the three
things the rerun actually changes: independent segment filtering, removal of
boundary-adjacent beats, and reducing the adaptation block from $500$ to $497$
labelled beats. Their magnitude is too small to alter any inferential
conclusion.

Third, and most importantly, \textbf{no inferential conclusion changes}.
Table~\ref{tab:splitfirst_contrasts} recomputes the four contrasts that matter
on the split-first cohort, reporting both raw and family-corrected $p$-values.
These four form their own sensitivity family and are \emph{not} pooled with the
six pre-specified primary comparisons. Both
encoder arms remain significantly above the frozen model on the split-first
cohort (under the locked estimator A4 improves $16/21$ patients and A5 $17/21$;
the arm \macrof{} in Table~\ref{tab:splitfirst} is seeds mean-averaged, ESC-4),
and the two comparisons against maximum
head-only replay remain non-significant, so the paper's central refusal to claim
superiority over A1 is unaffected. The one substantive change is that the
\emph{margin} over A1 shrinks --- A4$-$A1 falls from $+0.017$ to $+0.004$ and
A5$-$A1 from $+0.023$ to $+0.017$ --- which weakens the encoder-over-head
ordering slightly rather than strengthening it. We report that in the direction
it points: the split-first evidence makes the frontier reading \emph{more}
appropriate, not less.

This is a sensitivity analysis, not a replacement primary analysis. The headline
numbers throughout the paper remain those of the released pipeline, so that they
correspond to the pre-specified plan and to the published artifacts; the
split-first rerun is reported here as the check that the boundary dependency
does not carry the result. A fully causal streaming front end --- which would
also have to replace the zero-phase filter, not merely the record boundary ---
remains future work.

% GENERATED FILE -- do not edit by hand.
% Regenerate with: make tables  (scripts/make_tables.py)
\begin{table}[H]
\centering
\small
\caption{Boundary-leakage sensitivity (E17). The split-first rerun cuts the raw signal at the boundary, discards a $2$\,s guard each side, and filters the pieces independently, so no adaptation-side sample can depend on a test-side raw sample. \textbf{Every conclusion is preserved.}}
\label{tab:splitfirst}
\begin{tabular}{lrrr}
\toprule
Arm & Whole-record & Split-first & $\Delta$ \\
\midrule
A0 (frozen) & $0.666$ & $0.666$ & $-0.000$ \\
A1 (head $+$ max replay) & $0.786$ & $0.793$ & $+0.007$ \\
A4 (LoRA $r{=}1$) & $0.803$ & $0.796$ & $-0.006$ \\
A5 (LoRA $r{=}2$) & $0.809$ & $0.809$ & $+0.000$ \\
\bottomrule
\end{tabular}
\end{table}


% GENERATED FILE -- do not edit by hand.
% Regenerate with: make tables  (scripts/make_tables.py)
\begin{table}[H]
\centering
\small
\caption{Paired contrasts recomputed on the split-first (E17) cohort, $n = 21$ patients, five seeds averaged within patient. \textbf{These four sensitivity contrasts form their own family and are not pooled with the six pre-specified primary comparisons}; the Holm column corrects within these four only. Both frozen-model contrasts remain significant and neither encoder arm clears maximum head-only replay, matching the primary cohort.}
\label{tab:splitfirst_contrasts}
\setlength{\tabcolsep}{4pt}
\begin{tabular}{lrrcrrr}
\toprule
Contrast & Mean $\Delta$ & Med. $\Delta$ & 95\% CI & Raw $p$ & Holm $p$ & Impr./Wors. \\
\midrule
\textbf{A4$-$A0} & $+0.130$ & $+0.040$ & $[+0.050,+0.231]$ & $0.0012$ & \textbf{$0.0047$} & 17 / 4 \\
\textbf{A5$-$A0} & $+0.143$ & $+0.039$ & $[+0.060,+0.245]$ & $0.0012$ & \textbf{$0.0047$} & 17 / 4 \\
A4$-$A1 & $+0.004$ & $+0.001$ & $[-0.021,+0.027]$ & $0.3135$ & $0.4359$ & 14 / 6 \\
A5$-$A1 & $+0.017$ & $+0.001$ & $[-0.014,+0.047]$ & $0.2180$ & $0.4359$ & 14 / 6 \\
\bottomrule
\end{tabular}

\vspace{0.4em}
{\small ``Raw $p$'' is the unadjusted paired Wilcoxon signed-rank value; ``Holm $p$'' corrects across these four sensitivity contrasts. Intervals are patient bootstraps ($10{,}000$ resamples).}
\end{table}



\subsection{RQ5 --- Robustness to alternative class-reweighted source checkpoints}

\textbf{Takeaway:} the adaptation gain persists under two alternative
class-reweighted source-training objectives. This supports robustness to the
source \emph{objective}, but does \textbf{not} establish robustness to a
\emph{stronger unseen-patient source model}, which was never demonstrated
(secondary evidence).

A key robustness question is whether the adaptation gain merely corrects an
unusually weak S-class source. We train five architecture-identical source
variants on MIT--BIH DS1 with five seeds---cross-entropy, effective-number class
weighting, and focal loss $\gamma\in\{1,2,3\}$---and measure them under the same
protocol (Table~\ref{tab:e11_source_status}).

\paragraph{We do not call these checkpoints stronger, because on the reported
metric they are not.} Intra-patient DS1 S sensitivity is $0.881$ for plain
cross-entropy, $0.843$ for focal $\gamma{=}2$, and $0.776$ for effective-number
weighting: the class-reweighted objectives are \emph{below} the cross-entropy
baseline on exactly the axis that motivated the concern, and they also cost
overall \macrof{}. Calling them \emph{stronger source checkpoints} would not be
supported by these results, so we describe them for what they are ---
\textbf{alternative source checkpoints trained with class-reweighted
objectives} --- and treat RQ5 as a test of whether the adaptation result is
\emph{robust to the choice of source objective}, not as adaptation from a
stronger starting point.

One further caution: the S sensitivities above are measured on an intra-patient
DS1 split, whereas the frozen-source weakness that motivates adaptation is an
unseen-patient DS2 property. These are different domains and we do not compare
them as if they were the same. Establishing which checkpoint is genuinely
stronger on unseen patients would require evaluating every variant frozen on DS2,
which we have not done; that is the experiment RQ5 would need to make a true
stronger-source claim.

Training five variants is not the same as adapting from five, and we separate
the two explicitly: all five variants were \emph{trained} and measured for source
quality (Table~\ref{tab:e11_source_status}), while adaptation was re-run from
\textbf{two} of them --- effective-number weighting and focal $\gamma{=}2$, the
two that most change the minority-class behaviour the concern is about
(Table~\ref{tab:strong_adapt}). The gain survives: adapting from the
effective-number source lifts per-patient \macrof{} (N/S/V) from $0.458$ (frozen)
to $0.599$ (both head-only and encoder LoRA), and from the focal-$\gamma{2}$
source from $0.515$ to $0.642$ (head) and $0.644$ (encoder LoRA). On the
encoder-versus-head question these two runs say less than they appear to:
\emph{the mean encoder-minus-head direction remains nonnegative under two
alternative class-reweighted source objectives, but the margins are small
($+0.0001$ and $+0.0024$) and no uncertainty is reported}, so they are
consistent with the primary-cohort finding of a non-significant direct
comparison rather than independent support for an ordering.
Adaptation from the remaining three trained
variants (cross-entropy and focal $\gamma \in \{1,3\}$) is not reported here, so
we claim only that \emph{the ordering was re-evaluated from two
class-reweighted source checkpoints}, not from five. The three RR-architecture
and CNN variants require new model graphs and were never trained; they are
labeled \texttt{not\_implemented\_not\_run} in the ledger.

\begin{table}[H]
\centering
\small
\caption{E11 source-checkpoint ledger. All five variants share the architecture of
the shipped backbone and differ only in the training objective; each is trained on
MIT--BIH DS1 with five seeds and scored on a held-out \emph{intra-patient} DS1
split. The right-hand column separates the checkpoints we \emph{trained} from
those we also \emph{adapted from}, because the robustness claim covers only the
latter. \textbf{None of the reweighted objectives beats plain cross-entropy on S
sensitivity} ($0.881$), so we do not describe them as stronger sources; RQ5 tests
robustness to the source objective, not adaptation from a better starting point.
These intra-patient DS1 figures are also not comparable to unseen-patient DS2
behaviour.}
\label{tab:e11_source_status}
\begin{tabular}{llrrc}
\toprule
Source variant & Status & \macrof{} & S sens. & Adapted from? \\
\midrule
Cross-entropy (\texttt{original\_ce})   & measured, 5 seeds & \textbf{0.742} & \textbf{0.881} & no \\
Effective-number weighting              & measured, 5 seeds & 0.618 & 0.776 & \textbf{yes} \\
Focal loss $\gamma = 1$                 & measured, 5 seeds & 0.741 & 0.878 & no \\
Focal loss $\gamma = 2$                 & measured, 5 seeds & 0.730 & 0.843 & \textbf{yes} \\
Focal loss $\gamma = 3$                 & measured, 5 seeds & 0.725 & 0.780 & no \\
\midrule
RR-conditioned token embedding          & not implemented   & ---   & ---   & --- \\
RR token before attention               & not implemented   & ---   & ---   & --- \\
Lightweight inter-patient CNN           & not implemented   & ---   & ---   & --- \\
\bottomrule
\end{tabular}

\vspace{0.4em}
{\small The three unimplemented rows require new model graphs and were never
trained; they are recorded here rather than omitted so the ledger matches the
experiment plan. Three further pre-existing \texttt{beat\_classifier} checkpoints
exist in the repository but were measured on an intra-patient random split rather
than the DS1-only protocol, so they are not comparable and are excluded from this
table.}
\end{table}


\begin{table}[H]
\centering
\small
\caption{Adaptation from two E11 class-reweighted source checkpoints
(per-patient \macrof{} over N/S/V, DS2, five seeds). The adaptation gain persists
under a changed source objective, which is the robustness question RQ5 answers.
The $\Delta_{\mathrm{A4-A1}}$ column is nonnegative in both cases but the margins
are small and \textbf{no uncertainty is reported for them}; it must not be read
as evidence of an encoder-over-head ordering.}
\label{tab:strong_adapt}
\begin{tabular}{lrrrr}
\toprule
Source-training objective & Frozen & Head (A1) & Encoder LoRA (A4) & $\Delta_{\mathrm{A4-A1}}$ \\
\midrule
Effective-number weighting & 0.458 & 0.599 & 0.599 & $+0.0001$ \\
Focal loss $\gamma=2$       & 0.515 & 0.642 & 0.644 & $+0.0024$ \\
\bottomrule
\end{tabular}
\end{table}

\subsection{RQ6 --- Cross-database validation}
\label{sec:rq6}

\textbf{Takeaway:} per-subject adaptation transfers across databases --- INCART
($6$ subjects) and SVDB ($46$) both show positive gains, with encoder LoRA again
with a higher mean than head-only replay on both external subsets
(secondary evidence, small INCART cohort; no paired direct-comparison interval
is reported).

External validity is evaluated on INCART (11 recordings from six unique
subjects) and the MIT--BIH
Supraventricular Arrhythmia Database (SVDB, 46 records), with a documented
preprocessing manifest (lead, resampling, R-peak source, beat window, RR
normalization, AAMI mapping, exclusions). The same chronological,
target-informed protocol is applied per record. In contrast to a database-level
aggregate evaluation---which is neutral-to-negative because it performs no
per-record adaptation---the per-record result is clearly positive
(Table~\ref{tab:e13_crossdb}, Fig.~\ref{fig:crossdb}).

\paragraph{Statistical unit: unique subjects, not recordings.} INCART contains
multiple recordings per subject, so a record-level analysis would treat
non-independent recordings as independent samples. We parsed the subject
identifier from every INCART header and grouped accordingly: the 11 evaluated
records collapse to \textbf{six unique subjects} (\texttt{I01--I02} are one
65-year-old subject, \texttt{I03--I05} a second, \texttt{I06--I07} a third,
\texttt{I08} a fourth, \texttt{I09--I10} a fifth, and \texttt{I13} a sixth).
Because the saved cross-database cells retain the source record name, we
re-aggregated the existing runs to one score per subject---mean over that
subject's recordings and seeds---and bootstrap over subjects rather than
recordings. All INCART numbers below are subject-level. SVDB recordings are one
per subject, so its unit is unchanged. Every configured recording in both
databases now carries exactly one deterministic disposition
(Table~\ref{tab:external_manifest}).

\paragraph{What is adapted, and what is aggregated.} These two things are not the
same and the distinction matters. Adaptation is \textbf{record-wise}: each
recording starts from the same frozen source checkpoint, adapts on its own early
chronological segment, and is tested on its own later segment, with the model
\emph{reset between recordings}. Only the resulting scores are then averaged
within a subject, and the bootstrap resamples subjects. The correct description
is therefore \emph{record-wise chronological adaptation, aggregated to the
subject level for statistical analysis}. We do \textbf{not} carry one persistent
model across a subject's recordings, so nothing here should be read as evidence
about a single long-lived patient model that accumulates across sessions. True
subject-level adaptation---order a subject's recordings chronologically, adapt
once on the earliest, test on the later ones---is a different and stronger
experiment, and it is left to future work. Grouping to subjects fixes the
\emph{independence} defect in the statistics; it does not convert the protocol
into a longitudinal one.

\paragraph{The external gain survives the correction.} On INCART, subject-level
\macrof{} (N/S/V) rises from $0.459$ (frozen) to $0.595$ with head-only replay
($+0.137$, 95\% CI $[+0.058,+0.215]$), $0.618$ with encoder LoRA and raw replay
($+0.160$, CI $[+0.082,+0.234]$), and $0.647$ with no-replay encoder LoRA
($+0.189$, CI $[+0.110,+0.264]$); \textbf{6 of 6 subjects improve under every
method} and no interval includes zero. Collapsing recordings to subjects lowers
the encoder-LoRA gain from the record-level $+0.194$ to $+0.160$---the
record-level figure over-weighted the multi-recording subjects---but does not
change the conclusion. With $n = 6$ the inferential power is limited, so we treat
the INCART result as descriptive evidence with an interval attached rather than
as a significance claim.

On SVDB it rises from $0.384$ to $0.580$/$0.610$/$0.626$
($+0.196$/$+0.226$/$+0.242$), with 45--46 of 46 subjects improved and intervals
excluding zero (Table~\ref{tab:e13_crossdb}, Fig.~\ref{fig:crossdb}). Two
observations matter. First, the replay--plasticity ordering \emph{reproduces
externally and with a larger margin} than on DS2: encoder LoRA has a higher mean
than maximum head-only replay by $+0.023$ on INCART and $+0.029$ on SVDB. No
paired direct-comparison interval is reported for either external subset, so
these are mean differences and not a claim of statistical superiority. Second, dropping
replay entirely yields the \emph{largest} external target gain but also the most
source-domain change ($-0.11$ on INCART, $-0.14$ on SVDB, versus $\le -0.07$ with
replay), a textbook stability--plasticity trade-off: under strong domain shift,
replay trades a little target gain for markedly better source retention.

One caveat is not statistical but selective. The INCART cohort is a
tractability-capped convenience sample: the configuration attempts only the first
13 of 75 records, of which one fails preprocessing and one lacks minority-class
support. The 62 records outside that cap were never attempted, so the INCART
cohort is not a principled selection from the database and should not be read as
representative of it. The SVDB cohort is capped at 54 of 78 recordings on the
same basis, with 8 of the attempted 54 excluded for insufficient S$+$V support in
the test window.

\begin{table}[H]
\centering
\small
\caption{External-database inclusion and exclusion flow. \textbf{External
evaluation uses convenience-capped, minority-support-filtered subsets. The
eligibility rule uses test-window class labels} --- a recording is retained only
if its later test window carries at least 30 S$+$V beats --- \textbf{and
therefore enriches the cohort for recordings where minority-class adaptation can
be evaluated at all.} These are consequently not representative full-database
validations and no claim of representative external validation is made from them.
Every configured recording nevertheless carries exactly one deterministic
disposition; the full per-record manifests (lead, beat counts, class support, and
reason) are released with the code as
\texttt{incart\_inclusion\_exclusion\_manifest.csv} and
\texttt{svdb\_inclusion\_exclusion\_manifest.csv}.}
\label{tab:external_manifest}
\begin{tabular}{lrrrr}
\toprule
Stage & \multicolumn{2}{c}{INCART} & \multicolumn{2}{c}{SVDB} \\
\cmidrule(lr){2-3}\cmidrule(lr){4-5}
 & Records & Subjects & Records & Subjects \\
\midrule
Configured in the database        & 75 & 32 & 78 & 78 \\
Attempted (within record cap)     & 13 & 7  & 54 & 54 \\
Preprocessing successful          & 12 & 6  & 54 & 54 \\
Meets minority-class requirement  & 11 & 6  & 46 & 46 \\
\textbf{Final evaluated cohort}   & \textbf{11} & \textbf{6} & \textbf{46} & \textbf{46} \\
\midrule
\multicolumn{5}{l}{\emph{Exclusion reasons}} \\
Outside configured record cap     & 62 & --- & 24 & 24 \\
Preprocessing failure             &  1 & --- &  0 & 0 \\
Insufficient S$+$V test support   &  1 & --- &  8 & 8 \\
\bottomrule
\end{tabular}
\end{table}


\begin{table}[H]
\centering
\small
\caption{E13 cross-database adaptation, \textbf{aggregated to unique subjects}.
INCART recordings are grouped by the subject identifier in the WFDB header
(11 recordings $\rightarrow$ 6 subjects); SVDB recordings are one per subject.
Scores are per-subject \macrof{} over N/S/V with 95\% bootstrap intervals over
subjects, mean paired gain versus the frozen source, and subjects improved. With
$n = 6$ the INCART rows are descriptive rather than inferential.}
\label{tab:e13_crossdb}
\setlength{\tabcolsep}{3.5pt}
\begin{tabular}{llrcrcr}
\toprule
DB & Method & \macrof{} & 95\% CI & $\Delta$ & $\Delta$ CI & Impr. \\
\midrule
\multirow{4}{*}{\shortstack[l]{INCART\\$n=6$ subj.}}
 & Frozen source         & 0.459 & $[0.372,0.548]$ & ---      & ---                 & ---  \\
 & Head $+$ post-pool     & 0.595 & $[0.567,0.621]$ & $+0.137$ & $[+0.058,+0.215]$   & 6/6  \\
 & Enc.\ LoRA $+$ raw    & 0.618 & $[0.601,0.639]$ & $+0.160$ & $[+0.082,+0.234]$   & 6/6  \\
 & Enc.\ LoRA, no replay & 0.647 & $[0.630,0.666]$ & $+0.189$ & $[+0.110,+0.264]$   & 6/6  \\
\midrule
\multirow{4}{*}{\shortstack[l]{SVDB\\$n=46$ subj.}}
 & Frozen source         & 0.384 & $[0.335,0.434]$ & ---      & ---                 & ---   \\
 & Head $+$ post-pool     & 0.580 & $[0.543,0.618]$ & $+0.196$ & $[+0.152,+0.241]$   & 45/46 \\
 & Enc.\ LoRA $+$ raw    & 0.610 & $[0.566,0.653]$ & $+0.226$ & $[+0.179,+0.274]$   & 46/46 \\
 & Enc.\ LoRA, no replay & 0.626 & $[0.578,0.674]$ & $+0.242$ & $[+0.190,+0.297]$   & 45/46 \\
\bottomrule
\end{tabular}

\vspace{0.4em}
{\small Source-domain \macrof{} change on DS1 N/S/V: INCART
$-0.039$ / $-0.045$ / $-0.115$ and SVDB $-0.066$ / $-0.068$ / $-0.137$ for
head-only replay, encoder LoRA with raw replay, and encoder LoRA without replay
respectively. Dropping replay buys the largest target gain at the largest
retention cost.}
\end{table}


\begin{figure}[H]
\centering
\includegraphics[width=0.9\textwidth]{fig_crossdb_summary.pdf}
\caption{Subject-level cross-database adaptation on INCART and SVDB. INCART
recordings are grouped by the subject identifier in the WFDB header, giving six
unique subjects from 11 recordings; SVDB recordings are one per subject. Every
target-informed method improves per-subject \macrof{} (N/S/V) over the frozen
source (dashed line); annotations give the mean paired gain and the fraction of
subjects improved, and error bars are 95\% bootstrap intervals over subjects.
Encoder LoRA has a higher mean than head-only replay on both external subsets;
no paired direct-comparison interval is reported, so this is not a claim of
statistical superiority. With $n = 6$ the INCART
panel is descriptive rather than inferential.}
\label{fig:crossdb}
\end{figure}

\subsection{RQ7 --- Labeling burden and class coverage}
\label{sec:rq7}

\textbf{Takeaway:} around 100 labelled beats recover roughly 80\% of the
cohort-mean adaptation gain, but \emph{which} labels matter more than how many
--- coverage of the patient's minority class roughly doubles per-patient recovery
(secondary evidence).

Table~\ref{tab:e12_label_curve} and Fig.~\ref{fig:labels} give the label-budget
curve for A4 on the primary cohort under class-aware-oracle and repeated-random
selection at $K \in \{25,50,100,200,500\}$, with patient-bootstrap bands.
\textbf{The class-aware policy is a deliberately non-deployable oracle: it uses
knowledge of which classes later appear in the patient's test segment to choose
which beats to label, and therefore estimates an upper bound on the value of
class coverage rather than a policy a device could execute.} The random policy
is the deployable comparator; where the two differ, the gap is the headroom a
realisable selector would be competing for.
Measured on \emph{cohort means}, the class-aware policy recovers $0.58$ of the
500-label gain at 25 labels, $0.72$ at 50, $0.83$ at 100, and $0.93$ at 200;
random selection recovers $0.54$/$0.64$/$0.78$/$0.89$.

\paragraph{The cohort mean overstates what an individual patient gets.} The
cohort figure is a ratio of mean gains. The quantity a deployment actually cares
about is the per-patient recovery ratio
$R_K = (M_K - M_0)/(M_{500} - M_0)$, and averaging \emph{those} gives a
noticeably weaker picture (Table~\ref{tab:recovery}). Under class-aware
selection the median patient recovers $0.27$ at 25 labels, $0.61$ at 50, $0.67$
at 100, and $0.94$ at 200 --- and \textbf{only $38\%$ of patients individually
reach $R_K \ge 0.8$ at 100 labels}, rising to just $52\%$ at 200. The per-patient
\emph{mean} $R_K$ is unstable at small $K$ (it goes negative) because patients
whose 500-label gain is near zero have a near-zero denominator; we therefore
report the median and the fraction above threshold rather than the mean.

The defensible claim is consequently the narrow one: \emph{around 100 labelled
beats recover roughly 80\% of the cohort-mean adaptation gain under the tested
policies}. It is \textbf{not} true that 100 labels give 80\% of the gain for a
typical individual patient, and we do not claim that 50 labels suffice.

\paragraph{Which labels matter more than how many.} Splitting patients by
whether the adaptation window happened to contain the arrhythmia class that
later appears in the test window is the strongest signal in this experiment. At
100 class-aware labels, patients with that coverage recover $R_K = 0.82$ on
average against $0.41$ for those without --- roughly double. Label \emph{budget}
is therefore the wrong lever to optimize on its own: a small annotated set that
happens to include the patient's minority beats outperforms a larger one that
does not, which argues for coverage-aware annotation prompting rather than
simply asking a clinician for more labels.

\begin{table}[H]
\centering
\small
\caption{E12 label-budget curve for A4 on the 21-record patient-disjoint cohort
(5 seeds). ``Cohort recovery'' is the fraction of the 500-label mean gain over A0
recovered at that budget. Per-patient recovery, which is materially weaker, is in
Table~\ref{tab:recovery}. \emph{Values are seeds mean-averaged within patient;
median unification of this secondary curve is deferred (\texttt{BLOCKERS.md},
ESC-4).}}
\label{tab:e12_label_curve}
\begin{tabular}{lrrrr}
\toprule
Policy & Labels & \macrof{} & 95\% CI & Cohort recovery \\
\midrule
\multirow{5}{*}{Class-aware oracle}
 & 25  & 0.745 & $[0.672,0.818]$ & 0.582 \\
 & 50  & 0.764 & $[0.692,0.835]$ & 0.716 \\
 & 100 & 0.780 & $[0.705,0.853]$ & 0.832 \\
 & 200 & 0.793 & $[0.717,0.867]$ & 0.931 \\
 & 500 & 0.802 & $[0.733,0.873]$ & 1.000 \\
\midrule
\multirow{5}{*}{Repeated random}
 & 25  & 0.740 & $[0.669,0.813]$ & 0.541 \\
 & 50  & 0.754 & $[0.681,0.828]$ & 0.643 \\
 & 100 & 0.773 & $[0.698,0.848]$ & 0.784 \\
 & 200 & 0.787 & $[0.713,0.858]$ & 0.885 \\
 & 500 & 0.802 & $[0.731,0.875]$ & 1.000 \\
\midrule
\multicolumn{5}{l}{\emph{Reference points, full 500 labels, chronological}} \\
 A0 frozen & 500 & 0.666 & $[0.563,0.766]$ & --- \\
 A1 head $+$ max replay & 500 & 0.786 & $[0.712,0.859]$ & --- \\
 A5 LoRA $r{=}2$ & 500 & 0.809 & $[0.734,0.883]$ & --- \\
\bottomrule
\end{tabular}
\end{table}


\begin{table}[H]
\centering
\small
\caption{Label-budget recovery on the 21-record patient-disjoint cohort (5 seeds).
``Cohort'' is the ratio of mean gains, the figure usually quoted; $R_K$ is the
per-patient ratio $(M_K - M_0)/(M_{500} - M_0)$. The per-patient \emph{mean} is
unstable at small $K$ --- patients whose 500-label gain is near zero give a
near-zero denominator --- so the median and the fraction clearing $R_K \ge 0.8$
are the meaningful summaries. The last two columns split patients by whether the
adaptation window contained the arrhythmia class that later appears in the test
window.}
\label{tab:recovery}
\setlength{\tabcolsep}{5pt}
\begin{tabular}{lrrrrrrr}
\toprule
Policy & $K$ & Cohort & Mean $R_K$ & Median $R_K$ & $R_K \ge 0.8$ & $R_K$ covered & $R_K$ not \\
\midrule
\multirow{5}{*}{Class-aware oracle}
 & 25  & 0.582 & $-0.75$ & 0.27 & 19\% & $-0.47$ & $-0.94$ \\
 & 50  & 0.716 & $-0.03$ & 0.61 & 38\% & $0.31$  & $-0.25$ \\
 & 100 & 0.832 & $0.57$  & 0.67 & 38\% & \textbf{0.82} & 0.41 \\
 & 200 & 0.931 & $1.08$  & 0.94 & 52\% & $0.66$  & 1.37 \\
 & 500 & 1.000 & $1.00$  & 1.00 & 100\% & 1.00   & 1.00 \\
\midrule
\multirow{5}{*}{Repeated random}
 & 25  & 0.541 & $-0.56$ & 0.29 & 24\% & $-2.61$ & $-0.36$ \\
 & 50  & 0.643 & $0.02$  & 0.38 & 29\% & $0.17$  & $-0.02$ \\
 & 100 & 0.784 & $0.53$  & 0.66 & 29\% & 0.65    & 0.50 \\
 & 200 & 0.885 & $1.01$  & 0.84 & 52\% & 0.94    & 1.03 \\
 & 500 & 1.000 & $1.00$  & 1.00 & 100\% & 1.00   & 1.00 \\
\bottomrule
\end{tabular}
\end{table}


\begin{figure}[H]
\centering
\includegraphics[width=0.78\textwidth]{fig_label_budget_curve.pdf}
\caption{Label budget for A4 on the primary cohort, shown two ways because the
two disagree. \textbf{Left:} cohort-mean \macrof{} with patient-bootstrap 95\%
bands --- the curve usually reported, on which class-aware selection reaches
$\approx 80\%$ of the 500-label gain near 100 labels. \textbf{Right:} the
fraction of \emph{individual} patients whose recovery ratio $R_K$ clears $0.8$
(dotted line). That fraction is $0.38$ at 100 labels and still only $0.52$ at
200, so the left panel must not be read as a per-patient guarantee
(Table~\ref{tab:recovery}).}
\label{fig:labels}
\end{figure}

\subsection{Source-domain change}
\label{sec:srcchange}

\textbf{Takeaway:} tens of raw replay exemplars keep source-domain loss small and
bounded --- encoder arms do not incur larger source-domain loss than head-only
arms despite touching the encoder, and no source class collapses.

Across A1--A5 the source-domain \macrof{} change on the held-out DS1 N/S/V
classes stays within $[-0.048,-0.027]$ (frozen reference $0.946$), and the
five-class change stays within $\pm0.024$. Encoder arms do not incur larger
source-domain loss than head-only arms despite touching the encoder: tens of raw
replay exemplars are enough to anchor it, and no source class collapses. Fig.~\ref{fig:srcchange}
plots per-configuration target gain against source-domain change; the encoder arms
sit in the high-gain, small-source-loss region rather than trading one for the
other.

\begin{figure}[H]
\centering
\includegraphics[width=0.80\textwidth]{fig_gain_vs_forgetting.pdf}
\caption{Target patient \macrof{} gain versus the signed source-domain
\macrof{} change $\Delta^{\mathrm{src}}$ (Eq.~\ref{eq:gain_and_srcchange}) for
sweep configurations. $\Delta^{\mathrm{src}}$ is negative when adaptation costs
source performance, so points nearer the top-right retain more of the source
while gaining more on the target. Target gains are not bought with a
proportional source-domain loss; the encoder arms sit near the
stability--plasticity Pareto edge. The axis is labelled source-domain
\macrof{} change throughout; we avoid the word ``forgetting'' because the
quantity measured is a change in DS1 score, not a mechanism.}
\label{fig:srcchange}
\end{figure}
